
\subsection{Thermo-mechanical coupling and viscous dissipation in highly anisotropic plate-beams: an insect-scale wing-like structural benchmark}

The primary technical objectives of this second numerical example are:

\begin{enumerate}
    \item To assess the algorithmic robustness of the formulation under large-deformation flexural and torsional kinematics.
    \item To evaluate the performance of the anisotropic constitutive model, using an elongated span-to-width aspect ratio that amplifies the anisotropy-induced bending-torsion coupling and renders it clearly visible.
    \item To quantify the spatial non-uniformity of viscous self-heating induced by the anisotropic, torsional strain state.
    \item To verify the material stability of the anisotropic model under the combined bending-torsion state by evaluating directional projections of the acoustic tensor.
\end{enumerate}

The anisotropic response is introduced through a transverse-isotropy invariant term added to the isotropic hyperelastic energy, built on the structural tensor of a single fiber direction oriented at \qty{45}{\degree} in-plane. This off-axis alignment is what couples bending and torsion under normal actuation; the invariant formulation itself follows \cite{XX}.

As in the planar actuator, no mass or inertia is resolved: the formulation targets the constitutive response, not structural dynamics. This is unproblematic for the in-plane, equibiaxial actuation of the membrane, but it rules out a genuine multi-cycle actuation here, since a compliant, elongated cantilever driven out-of-plane over many periods would be expected to show inertial effects that a quasi-static solve cannot capture. The temporal, multi-cycle analysis of self-heating is therefore reported for the membrane benchmark only; this example instead reports the single actuation snapshot below, and focuses on the spatial structure of the anisotropy-induced self-heating that the homogeneous membrane state cannot exhibit.

\begin{figure}
  \centering
  \includegraphics[width=.7\textwidth]{beam/beam_geometry.png}
  \caption{Discrete model of the plate beam. The black lines represent the active and passive layers of electrodes triggering the bending.}
  \label{fig:beam_geom}
\end{figure}


\begin{figure}[htbp]
  \centering
  \begin{subfigure}[t]{.49\textwidth}
    \centering
    \includegraphics[width=\textwidth]{beam/deformed_mesh.png}
    \caption{Deformed mesh.}
  \end{subfigure}
  \begin{subfigure}[t]{.49\textwidth}
    \centering
    \includegraphics[width=\textwidth]{beam/displacement_magnitude.png}
    \caption{Displacement magnitude.}
  \end{subfigure}
  \\
  \begin{subfigure}[t]{.49\textwidth}
    \centering
    \includegraphics[width=\textwidth]{beam/electric_potential.png}
    \caption{Electric potential.}
  \end{subfigure}
  \begin{subfigure}[t]{.49\textwidth}
    \centering
    \includegraphics[width=\textwidth]{beam/temperature.png}
    \caption{Temperature, localized by the thermoelastic coupled state.}
  \end{subfigure}
  \\
  \begin{subfigure}[t]{.49\textwidth}
    \centering
    \includegraphics[width=\textwidth]{beam/first_piola_11.png}
    \caption{First Piola-Kirchoff $P_{11}$.}
  \end{subfigure}
  \begin{subfigure}[t]{.49\textwidth}
    \centering
    \includegraphics[width=\textwidth]{beam/first_piola_22.png}
    \caption{First Piola-Kirchoff $P_{22}$.}
  \end{subfigure}
  \\
  \begin{subfigure}[t]{.49\textwidth}
    \centering
    \includegraphics[width=\textwidth]{beam/first_piola_33.png}
    \caption{First Piola-Kirchoff $P_{33}$.}
  \end{subfigure}
  \begin{subfigure}[t]{.49\textwidth}
    \centering
    \includegraphics[width=\textwidth]{beam/first_piola_23.png}
    \caption{First Piola-Kirchoff $P_{23}$.}
  \end{subfigure}
  \\
  % \begin{subfigure}[t]{.49\textwidth}
  %   \centering
  %   \includegraphics[width=\textwidth]{beam/first_piola_13.png}
  %   \caption{First Piola-Kirchoff $P_{13}$.}
  % \end{subfigure}
  % \begin{subfigure}[t]{.49\textwidth}
  %   \centering
  %   \includegraphics[width=\textwidth]{beam/first_piola_23.png}
  %   \caption{First Piola-Kirchoff $P_{23}$.}
  % \end{subfigure}
  \caption{Simulation results of the plate beam at time $t=\qty{0.025}{second}$ during the maximum amplitude of the first cycle of the electric potential.}
  \label{fig:beam_screenshots}
\end{figure}


% \subsection{Conservation of properties}

\begin{figure}[htbp]
  \centering
  \includegraphics[width=.9\textwidth]{beam/time_evolution_properties.png}
  \caption{Time evolution during $0.5s$. (\emph{Top}) Pitch and stroke at the tip of the beam; \emph{Center} Viscous dissipation; \emph{Bottom} Average temperature.}
  \label{fig:time_evolution}
\end{figure}


% \subsubsection{Material stability: acoustic tensor projections}

% The acoustic tensor is evaluated over the full range of propagation directions $(\alpha,\beta)$ at the state of maximum bending-torsion coupling. Strong ellipticity is retained under this combined flexural-torsional, anisotropic loading, a stricter test than the pre-stretched state checked in the planar actuator example.

% \figplaceholder{7cm}{Directional projection of the acoustic tensor (minimum eigenvalue) over $(\alpha,\beta) \in [0,2\pi]\times[0,\pi]$, evaluated at the peak bending-torsion state of the elongated beam. Source: \texttt{examples/AcousticTensor/AcousticTensor.jl}.}
% % \label{fig:beam_acoustic_tensor}


